% <<<PDFPAGE 214>>>
\begin{chapterreferences}
\refentry{1} Yann LeCun, Sumit Chopra, Raia Hadsell, M Ranzato, and F Huang. A tutorial on energy-based learning. \emph{Predicting structured data}, 1(0), 2006.
\refentry{2} David H Ackley, Geoffrey E Hinton, and Terrence J Sejnowski. A learning algorithm for Boltzmann machines. \emph{Cognitive science}, 9(1):147–169, 1985.
\refentry{3} Paul Smolensky. \emph{Information processing in dynamical systems: Foundations of harmony theory}. Technical report, Colorado Univ at Boulder Dept of Computer Science, 1986.
\refentry{4} Edwin T Jaynes. \emph{Probability theory: The logic of science}. Cambridge university press, 2003.
\refentry{5} Will Grathwohl, Kuan-Chieh Wang, Joern-Henrik Jacobsen, David Duvenaud, Mohammad Norouzi, and Kevin Swersky. Your classifier is secretly an energy-based model and you should treat it like one. In \emph{International Conference on Learning Representations}, 2019.
\refentry{6} Max Welling and Yee W Teh. Bayesian learning via stochastic gradient Langevin dynamics. In \emph{Proceedings of the 28th international conference on machine learning (ICML-11)}, pages 681–688. Citeseer, 2011.
\refentry{7} Geoffrey E Hinton, Terrence J Sejnowski, et al. Learning and relearning in Boltzmann machines. \emph{Parallel distributed processing: Explorations in the microstructure of cognition}, 1(282–317):2, 1986.
\refentry{8} Benjamin Marlin, Kevin Swersky, Bo Chen, and Nando Freitas. Inductive principles for restricted Boltzmann machine learning. In \emph{Proceedings of the thirteenth International Conference on Artificial Intelligence and Statistics}, pages 509–516, 2010.
\refentry{9} Geoffrey E Hinton. Training products of experts by minimizing contrastive divergence. \emph{Neural computation}, 14(8):1771–1800, 2002.
\refentry{10} Tijmen Tieleman. Training restricted Boltzmann machines using approximations to the likelihood gradient. In \emph{ICML}, pages 1064–1071, 2008.
\refentry{11} Jakub M Tomczak and Adam Gonczarek. Learning invariant features using subspace restricted Boltzmann machine. \emph{Neural Processing Letters}, 45(1):173–182, 2017.
\refentry{12} Roland Memisevic and Geoffrey E Hinton. Learning to represent spatial transformations with factored higher-order Boltzmann machines. \emph{Neural computation}, 22(6):1473–1492, 2010.
% <<<PDFPAGE 215>>>
\refentry{13} Aaron Courville, James Bergstra, and Yoshua Bengio. A spike and slab restricted Boltzmann machine. In \emph{Proceedings of the fourteenth international conference on artificial intelligence and statistics}, pages 233–241. JMLR Workshop and Conference Proceedings, 2011.
\refentry{14} Ruslan Salakhutdinov, Andriy Mnih, and Geoffrey Hinton. Restricted Boltzmann machines for collaborative filtering. In \emph{Proceedings of the 24th international conference on Machine learning}, pages 791–798, 2007.
\refentry{15} Kyung Hyun Cho, Tapani Raiko, and Alexander Ilin. Gaussian-Bernoulli deep Boltzmann machine. In \emph{The 2013 International Joint Conference on Neural Networks (IJCNN)}, pages 1–7. IEEE, 2013.
\refentry{16} Ilya Sutskever, Geoffrey E Hinton, and Graham W Taylor. The recurrent temporal restricted Boltzmann machine. In \emph{Advances in Neural Information Processing Systems}, pages 1601–1608, 2009.
\refentry{17} Graham W Taylor, Leonid Sigal, David J Fleet, and Geoffrey E Hinton. Dynamical binary latent variable models for 3d human pose tracking. In \emph{2010 IEEE Computer Society Conference on Computer Vision and Pattern Recognition}, pages 631–638. IEEE, 2010.
\refentry{18} Hugo Larochelle and Yoshua Bengio. Classification using discriminative restricted Boltzmann machines. In \emph{Proceedings of the 25th international conference on Machine learning}, pages 536–543, 2008.
\refentry{19} Hugo Larochelle, Michael Mandel, Razvan Pascanu, and Yoshua Bengio. Learning algorithms for the classification restricted Boltzmann machine. \emph{The Journal of Machine Learning Research}, 13(1):643–669, 2012.
\refentry{20} Jakub M Tomczak. Learning informative features from restricted Boltzmann machines. \emph{Neural Processing Letters}, 44(3):735–750, 2016.
\refentry{21} Jakub M Tomczak. On some properties of the low-dimensional Gumbel perturbations in the Perturb-and-MAP model. \emph{Statistics \& Probability Letters}, 115:8–15, 2016.
\refentry{22} Jakub M Tomczak, Szymon Zaręba, Siamak Ravanbakhsh, and Russell Greiner. Low-dimensional perturb-and-map approach for learning restricted Boltzmann machines. \emph{Neural Processing Letters}, 50(2):1401–1419, 2019.
\refentry{23} Jascha Sohl-Dickstein, Peter B Battaglino, and Michael R DeWeese. New method for parameter estimation in probabilistic models: Minimum Probability Flow. \emph{Physical review letters}, 107(22):220601, 2011.
\refentry{24} Yang Song and Diederik P Kingma. How to train your energy-based models. \emph{arXiv preprint arXiv:2101.03288}, 2021.
\refentry{25} Yoshua Bengio. \emph{Learning deep architectures for AI}. Now Publishers Inc, 2009.
\refentry{26} Honglak Lee, Roger Grosse, Rajesh Ranganath, and Andrew Y Ng. Convolutional deep belief networks for scalable unsupervised learning of hierarchical representations. In \emph{Proceedings of the 26th annual international conference on machine learning}, pages 609–616, 2009.
\refentry{27} Ruslan Salakhutdinov. Learning deep generative models. \emph{Annual Review of Statistics and Its Application}, 2:361–385, 2015.
% <<<PDFPAGE 216>>>
\refentry{28} Ruslan Salakhutdinov and Iain Murray. On the quantitative analysis of deep belief networks. In \emph{Proceedings of the 25th international conference on Machine learning}, pages 872–879, 2008.
\refentry{29} Geoffrey E Hinton and Ruslan R Salakhutdinov. Reducing the dimensionality of data with neural networks. \emph{science}, 313(5786):504–507, 2006.
\refentry{30} Max Welling and Yee Whye Teh. Approximate inference in Boltzmann machines. \emph{Artificial Intelligence}, 143(1):19–50, 2003.
\refentry{31} Jonathan S Yedidia, William T Freeman, and Yair Weiss. Constructing free-energy approximations and generalized belief propagation algorithms. \emph{IEEE Transactions on information theory}, 51(7):2282–2312, 2005.
\refentry{32} Martin J Wainwright, Tommi S Jaakkola, and Alan S Willsky. A new class of upper bounds on the log partition function. \emph{IEEE Transactions on Information Theory}, 51(7):2313–2335, 2005.
\refentry{33} Tamir Hazan and Tommi Jaakkola. On the partition function and random maximum a-posteriori perturbations. In \emph{Proceedings of the 29th International Conference on International Conference on Machine Learning}, pages 1667–1674, 2012.
\end{chapterreferences}
